% =============================================================================
%  Brute-Force Solver
% =============================================================================
% Requires in the parent document:
% \usepackage{algorithm}
% \usepackage{algpseudocode}

\section{Brute-Force Solver}
\label{sec:brute_force}

\subsection{Role in the Project}

\texttt{BruteForceSolver} is the simplest solver in the repository. It does
not try to be clever. It fills every empty cell with every possible value,
checks the completed grid, and returns the first grid that satisfies the
Futoshiki rules.

This makes it a useful baseline. The code is easy to inspect, and its behaviour
is predictable on small puzzles. It is not a serious option for large
instances. Once the number of empty cells grows, the candidate count grows too
quickly.

\subsection{Public Interface}

The class uses the same \texttt{BaseSolver} contract as the other solvers:

\texttt{solve(puzzle, on\_step=None)} returns
\texttt{(Puzzle | None, Stats)}, and \texttt{get\_name()} returns a string.

\texttt{get\_name()} returns \texttt{"BruteForce"}. \texttt{solve()} receives a
\texttt{Puzzle} and returns either a solved \texttt{Puzzle} or \texttt{None},
plus a \texttt{Stats} object. The input puzzle is left alone. Each candidate is
written into a copied puzzle.

The optional \texttt{on\_step} callback is used by the UI. In this solver it is
called as \texttt{on\_step(candidate.grid.copy())}, so the UI receives a grid
snapshot rather than the mutable candidate puzzle.

If the callback raises \texttt{StopIteration}, the solver stops cleanly and
returns the statistics collected so far.

\subsection{Search Space}

Let $N$ be the grid size and let $K$ be the number of empty cells in the input
puzzle. Each empty cell can receive any value from $1$ to $N$, so the candidate
space is

\begin{equation}
  |\mathcal{C}| = N^{K}.
  \label{eq:bf_candidates}
\end{equation}

For a fully empty $N \times N$ grid, $K = N^{2}$.

\begin{table}[h]
  \centering
  \caption{Candidate space size for fully empty grids.}
  \label{tab:bf_space}
  \begin{tabular}{ccr}
    \toprule
    $N$ & $K = N^{2}$ & $|\mathcal{C}| = N^{K}$ \\
    \midrule
    4 & 16 & $4^{16} \approx 4.3 \times 10^{9}$ \\
    5 & 25 & $5^{25} \approx 2.98 \times 10^{17}$ \\
    6 & 36 & $6^{36} \approx 1.50 \times 10^{28}$ \\
    7 & 49 & $7^{49} \approx 2.56 \times 10^{41}$ \\
    \bottomrule
  \end{tabular}
\end{table}

This is why the solver only makes sense for small or heavily filled puzzles. It
does not shrink a cell's domain before trying values. It simply walks through
the Cartesian product of all possible values for all empty cells.

\subsection{Implementation}

The code is a direct exhaustive search. It asks the \texttt{Puzzle} for its
empty cells, then treats each assignment as a complete candidate grid. The real
code uses \texttt{itertools.product} for the assignment loop. The same flow is
easier to read as pseudocode:

\begin{algorithm}[h]
  \caption{\textsc{BruteForceSolve}}
  \label{alg:brute_force_solve}
  \begin{algorithmic}[1]
    \Procedure{BruteForceSolve}{$puzzle, on\_step$}
      \State start memory and time tracking
      \State $initially\_unsolved \gets$ number of zeros in $puzzle.grid$
      \State $empty\_cells \gets \Call{GetEmptyCells}{puzzle}$
      \State $domain \gets \{1,2,\ldots,puzzle.N\}$
      \State $expansions \gets 0$
      \State $result \gets \textsc{None}$
      \ForAll{$assignment \in \Call{CartesianProduct}{domain, |empty\_cells|}$}
        \State $expansions \gets expansions + 1$
        \State $candidate \gets \Call{Copy}{puzzle}$
        \ForAll{$(cell, value) \in \Call{Zip}{empty\_cells, assignment}$}
          \State write $value$ into $candidate.grid[cell]$
        \EndFor
        \If{$on\_step$ exists}
          \State call $on\_step$ with a copy of $candidate.grid$
        \EndIf
        \If{\Call{IsValid}{$candidate$}}
          \State $result \gets candidate$
          \State \textbf{break}
        \EndIf
      \EndFor
      \State $stats \gets$ time, memory, $expansions$, and completion ratio
      \State \Return $(result, stats)$
    \EndProcedure
  \end{algorithmic}
\end{algorithm}

The real implementation wraps the loop in a small \texttt{try}/\texttt{except}
block. If the UI callback raises \texttt{StopIteration}, the loop stops and the
solver returns with the statistics collected up to that point.

This pseudocode hides Python details such as \texttt{time.perf\_counter()},
\texttt{tracemalloc}, and the exact \texttt{Stats} constructor, but it matches
the structure of \texttt{BruteForceSolver.solve()}.

A few details are worth calling out:

\begin{itemize}
  \item \texttt{itertools.product} is lazy, so the solver does not store all
        candidates at once.
  \item \texttt{expansions} is incremented once for each candidate assignment
        tried.
  \item \texttt{puzzle.copy()} creates a new \texttt{Puzzle} with an independent
        NumPy grid. The constraint objects are copied as list entries but still
        represent the same fixed puzzle constraints.
  \item The callback sees a grid snapshot before the candidate is validated.
\end{itemize}

There is no partial state object and no recursive call stack. The loop is flat:
generate one complete assignment, build one candidate puzzle, validate it, and
move on.

\subsubsection{Candidate Validation}

\texttt{\_is\_valid()} checks only completed candidates. It applies the same
three Futoshiki requirements every time:

\begin{enumerate}
  \item every row must be a permutation of $1..N$;
  \item every column must be a permutation of $1..N$;
  \item every horizontal and vertical inequality constraint must be satisfied.
\end{enumerate}

The row and column checks use sorted values:
each row and column must sort to \([1,2,\ldots,N]\).

Inequality checks are delegated to \texttt{InequalityConstraint.is\_satisfied()}.
It returns \texttt{False} as soon as one check fails.

In pseudocode, the helper is:

\begin{algorithm}[h]
  \caption{\textsc{IsValid}}
  \label{alg:brute_force_is_valid}
  \begin{algorithmic}[1]
    \Procedure{IsValid}{$candidate$}
      \State $expected \gets [1,2,\ldots,candidate.N]$
      \ForAll{$row \in candidate.grid$}
        \If{$\Call{Sorted}{row} \ne expected$}
          \State \Return \textsc{False}
        \EndIf
      \EndFor
      \ForAll{$column \in candidate.grid$}
        \If{$\Call{Sorted}{column} \ne expected$}
          \State \Return \textsc{False}
        \EndIf
      \EndFor
      \ForAll{constraint in candidate's horizontal and vertical constraints}
        \If{$constraint$ is not satisfied by $candidate$}
          \State \Return \textsc{False}
        \EndIf
      \EndFor
      \State \Return \textsc{True}
    \EndProcedure
  \end{algorithmic}
\end{algorithm}

Bad partial assignments are not pruned. If the first two values in a row
already duplicate each other, the solver still finishes the full candidate
tuple before rejecting it. That is the main difference between this solver and
the backtracking-based solvers in the project.

\subsection{Statistics}

\texttt{solve()} starts \texttt{tracemalloc}, records the start time with
\texttt{time.perf\_counter()}, and builds a \texttt{Stats} object before
returning.

\begin{table}[h]
  \centering
  \caption{Statistics reported by \texttt{BruteForceSolver}.}
  \label{tab:bf_stats}
  \begin{tabular}{lp{0.62\linewidth}}
    \toprule
    Field & Meaning in this solver \\
    \midrule
    \texttt{time\_ms} & wall-clock solving time in milliseconds \\
    \texttt{memory\_kb} & peak heap memory reported by \texttt{tracemalloc} \\
    \texttt{inference\_count} & always 0; no inference rules are used \\
    \texttt{node\_expansions} & number of complete candidate assignments tried \\
    \texttt{backtracks} & always 0; the solver does not backtrack \\
    \texttt{completion\_ratio} & \texttt{1.0} if solved, \texttt{0.0} if unsolved, with complete inputs treated as \texttt{1.0} \\
    \bottomrule
  \end{tabular}
\end{table}

The completion ratio comes from the number of cells that were empty at the
start. If the puzzle has no empty cells, the ratio is \texttt{1.0}. If no
solution is found, it is \texttt{0.0}. If a solution is found, all initially
empty cells have been filled, so the ratio is \texttt{1.0}.

\subsection{Complexity}

The worst-case number of candidates is $N^K$. For each candidate, the solver
copies a grid, fills $K$ cells, checks rows and columns, and checks all
inequality constraints. A useful summary is:

\begin{equation}
  T = O\!\left(N^{K} \cdot (K + N^{2}\log N + |\mathcal{I}|)\right),
\end{equation}

where $|\mathcal{I}|$ is the number of inequality constraints. Since
$K \le N^2$ and $|\mathcal{I}| \le 2N(N-1)$, the exponential term dominates.

Memory use is much smaller than the running time might suggest.
\texttt{itertools.product} keeps only its current tuple and cursor state. The
solver also keeps only one candidate puzzle at a time.

\begin{table}[h]
  \centering
  \caption{Memory used by the brute-force loop.}
  \label{tab:bf_space2}
  \begin{tabular}{ll}
    \toprule
    Component & Size \\
    \midrule
    Empty-cell list & $O(K)$ \\
    Current assignment tuple & $O(K)$ \\
    Candidate grid & $O(N^{2})$ \\
    Call stack & $O(1)$ \\
    \midrule
    \textbf{Total} & $O(N^{2})$ because $K \le N^{2}$ \\
    \bottomrule
  \end{tabular}
\end{table}

\subsection{Behaviour on Edge Cases}

The tests cover the cases most likely to break this style of solver:

\begin{itemize}
  \item a complete valid puzzle is returned after one expansion;
  \item a complete invalid puzzle returns \texttt{None};
  \item a puzzle with one empty cell tries at most $N$ values;
  \item row, column, and inequality constraints are all checked on the result;
  \item contradictory givens or impossible inequalities return \texttt{None};
  \item \texttt{inference\_count} and \texttt{backtracks} stay at 0.
\end{itemize}

These cases match the actual design: exhaustive, flat, and easy to reason
about. The price is that failed partial assignments teach the solver nothing.

\subsection{Design Tradeoffs}

\paragraph{Flat enumeration instead of recursion.}
The solver uses \texttt{itertools.product} rather than a recursive assignment
function. The code is shorter, the call stack stays flat, and CPython handles
the counter logic. The downside is that the solver has no natural place to
prune partial assignments.

\paragraph{Copy per candidate.}
Each candidate uses a fresh \texttt{Puzzle} copy. That avoids undo logic and
makes the loop easy to follow. It also allocates a new grid for every candidate,
which becomes wasteful once the search space grows.

\paragraph{Full validation after filling the grid.}
Validation happens only after all empty cells have values. This makes
\texttt{\_is\_valid()} simple and keeps the solver independent from partial
constraint logic. It also means the solver spends time completing candidates
that could have been rejected much earlier.

\paragraph{Baseline value.}
The brute-force solver is useful because it is plain. It gives the repository a
simple implementation of the solver interface and a clear picture of what
unpruned search looks like. Its limitations are part of the point: the other
solvers exist because this approach does not scale.
